\subsection{A reusable state-borrowing rule}

The concrete reduction is an instance of the following local rule.  Suppose
an old machine has distinct states $d,r,s,t,u,v$ with rows
\[
\delta(d,0)=(0,R,u),\qquad\delta(d,1)=(1,R,v),
\]
and
\[
\delta(r,0)=(0,L,s),\qquad
\delta(s,0)=(0,R,t),\qquad
\delta(t,0)=(0,R,u).
\]
Assume all old occurrences of $r$ read zero, as established by an inductive
invariant.  Assume every incoming row to $d$ is of the form
$(p,0)\mapsto(0,R,d)$, with none of $d,r,s,t$ used as such an entry state.
Finally require that the three displayed return rows remain present and
unchanged after deleting $d$.

Delete $d$, redirect each of its incoming rows to $r$, and replace the
previously inadmissible row $r/1$ by $(1,R,v)$.  The machine loses one
working state.  A formerly $d$-entering row writes zero immediately left of
the new $r$ configuration.  If $r$ reads one, its replacement row performs
the old $d/1$ action.  If $r$ reads zero, the three-step return
$r,s,t,u$ performs the old $d/0$ action with exactly two extra transitions.
The unchanged old occurrences of $r$ read zero and retain their old row.
The same strictly positive block induction as above proves halting
equivalence, provided the other ordinary-row checkpoint conditions hold.

For the present construction the substitution is
\[
(d,r,s,t,u,v,p)=(15,27,32,45,21,22,11).
\]
The required return is already in the decoder; it is not a newly allocated
subroutine.  The two symbols together with the earlier written tape cell
provide the discrimination that would otherwise require state $15$.

This lemma separates the reusable optimization from the finite proof that
a particular target branch was inadmissible.  Such a branch need not remain
inadmissible in the optimized machine: here $27/1$ becomes essential.  A
coverage report showing an unused row is therefore a useful search lead,
but the induction invariant and return identity are the actual proof.
